\documentclass[../../../include/open-logic-section]{subfiles}

\begin{document}
\olfileid{sth}{ord-arithmetic}{intro}

\olsection{પરિચય}

\olref[ordinals][]{chap}માં આપણે ક્રમવાચક સંખ્યાઓનો
સિદ્ધાંત વિકસાવ્યો. \olref[spine][]{chap}માં જોયું કે
ક્રમવાચક સંખ્યાઓને એવી ધરી માની શકીએ જેની આસપાસ
પદાનુક્રમનો બાકીનો ભાગ રચાય છે. પરંતુ તેમનું કામ માત્ર
એટલું જ નથી. ક્રમવાચક સંખ્યાઓનું અંકગણિત પણ કરવાનું છે.

\olref[ordinals][idea]{sec}માં $\omega$,
$\omega+1$ અને $\omega+\omega$ની વાત કરતાં આપણે
તેનો સંકેત આપી દીધો હતો. ત્યારે વાત અનૌપચારિક હતી;
હવે તેને ચોક્કસ રીતે સમજાવવાનો સમય છે. જોકે ઉલ્લેખ
કરવો જોઈએ કે આ પ્રકરણમાં તત્ત્વચિંતન બહુ ઓછું છે:
તકનિકી વિકાસ છે અને સાથે એક (થોડું) રસપ્રદ અવલોકન
છે કે એ જ કામ બે જુદી રીતે કરી શકીએ છીએ.

\end{document}
